In this small CPU study the answer depended on the training budget and on one data property. At the planned 400 steps, no nonlinear model passed the pre-registered rule against a DLinear-style model on any of six tasks at horizons 24, 96 or 192. At 1000 and 2000 steps, with DLinear's MSE essentially unchanged, the patch GRU passed it on regime-switching data at horizons 24 and 96 (seed-mean MSE \rhval{sweep_steps/gru@steps=2000/regime/mse_96/mean:3} against \rhval{sweep_steps/dlinear@steps=2000/regime/mse_96/mean:3} at $H=96$, 2000 steps) and the patch transformer at horizon 24; neither passed it at horizon 192. Regime switching was the only property we varied under which a nonlinear model passed the rule: on the four tasks without regime changes (tested at $H=96$ with 2000 steps) and on ETTh1 none did, and on ETTh1 DLinear remained better at $H=96$ and 192 at every budget. The GRU's wins were confined to mean regime lengths up to 500 steps and its margin did not grow monotonically with the switching rate, so H2 is only partly supported. Whether this holds with tuned or larger models, more seeds or longer budgets remains open.
